\scanpage{142}
% archive foliation: 141
la situation par $G_1$, $G'_1$, $(P_1,\xi_1)$, $(Q_1,\eta_1)$ etc.

[Attention, on perd le fait que $\unsure{S}$ est réduit, $S_1$ ne l'est plus].
$P_1$ s'\uline{identifie} alors à $G_1$
\inserted{(muni du pt marqué $\varepsilon_1$)},
$Q_1$ à $G'_1\times_{S_1}G'_1$ muni de $(\varepsilon'_1,\varepsilon'_1)$,
\deleted{\unsure{je dis que}} enfin $(P'_1,\xi'_1)$ s'identifie aussi à
$(G'_1,\varepsilon'_1)$.
Il suffit de voir que $(P'_1,\xi'_1)\to(Q_1,\eta_1)$ est
l'application diagonale $G'_1-G'_1\times_{S_1}G'_1$.

\uline{On} \unsure{fait} ces identifications, et désigne par $w$ l'isom.
$(P,\xi)\times_u G'-(P,\xi)\times_v G'$,
$w_1$ l'isom. déduit par extension de la base à $S_1$,
notre assertion signifie que $w_1$ transforme la section
\unsure{``naturelle''} de $P_1\times_{u_1}G'_1$ dans celle de
$P_1\times_{v_1}G'_1$, i.e. qu'on a commutativité dans

\diagram{\begin{tikzcd}[column sep=large]
P=S_1 \arrow[r] \arrow[rr,bend right=25] &
(P\times_u G')\times_S S_1
\arrow[r,"w_1=w\times_S\unsure{S'_1}"] &
(P\times_v G')\times_S S_1
\end{tikzcd}}

ou encore qu'on a commutativité de

\diagram{\begin{tikzcd}[column sep=large]
P \arrow[r,"\mathrm{can}"] \arrow[rr,bend right=25,"\mathrm{can}"'] &
P\times_u G' \arrow[r,"w"] & P\times_v G'
\end{tikzcd}}

\uline{Or cela résulte} de l'\uline{unicité} de $P\to P\times_vG'$
transformant \deleted{\unsure{$\xi$ en $\xi$}} pt marqué en pt marqué \unsure{...}.
